% =====================================================================
% P1 Appendix A.2 --- Proposition 1: value-error decomposition and
% the limits of uniform W_1 / Le Cam bounds.
%
% Source of truth: drafts/proofs-P3-P5-P6.md v3.1, Sections 3.1--3.4.5
% (including the F6 analysis and the three-term decomposition of 2026-10-20).
% This file is a fragment: it is meant to be \input under
%   \appendix \section{Proofs}
% after A.1 and before A.3.  The standalone driver main-a12.tex sets
% the subsection counter to 0, so this fragment becomes A.2.
%
% Requirements: amsmath, amssymb, amsthm, and hyperref (optional, but
% the fixed-tag mechanism below keeps destinations unique when it is
% loaded).  Equation tags (3.4.1)--(3.4.4) follow the supplementary development; do not
% renumber without updating the main text.
% =====================================================================

\allowdisplaybreaks

\providecommand{\TV}{\mathrm{TV}}

\newtheorem*{aIIPropMain}{Proposition 1}
\newtheorem*{aIIPropTV}{Proposition A.1}
\newtheorem*{aIILemOneStep}{Lemma A.2}
\newtheorem*{aIICorMixed}{Corollary A.3}
\newtheorem*{aIIDefLeCam}{Definition A.4}
\newtheorem*{aIILemLeCam}{Lemma A.5}
\newtheorem*{aIIPropLeCam}{Proposition A.6}
\newtheorem*{aIIPropWone}{Proposition A.7}
\newtheorem*{aIIPropBoundary}{Proposition A.8}

% Fixed equation tags do not advance the equation counter, so hyperref
% would give every display the same destination.  A dedicated counter
% keeps the destinations unique without changing the printed (frozen)
% tag numbers.  The original \theHequation is restored at the end.
\providecommand{\theHequation}{}
\let\aIItheHequation\theHequation
\newcounter{aIIKanchor}
\renewcommand{\theHequation}{aII.\arabic{aIIKanchor}}
\newcommand{\aIIfixedtag}[1]{\stepcounter{aIIKanchor}\tag{#1}}

\subsection{Value-error decomposition and the limits of uniform bounds (Proposition~1)}\label{app:a2}

This appendix formalizes Proposition~1 (three-term decomposition).  It is a negative result plus a
conditional decomposition, not an identification theorem.  It explains why the
paper does not report a single $W_1$ or Le Cam value-error bound for a fixed
policy, and what it reports instead.  The formal guarantees of the paper remain
Theorem~1 (population identification and plug-in consistency) and Theorem~2
(finite-sample coverage).  The value-error track is an estimable sensitivity
analysis, not a main guarantee.

\paragraph{Setting and notation.}
Consider two POMDPs that share $(X,A,T,\rho,\nu,\pi)$ and differ only in the
observation kernel, $h$ versus $h'$.  Write $V_h^\pi$ for the value under
kernel $h$.  For a belief $b$ and kernel $h$, define the predictive law and the
posterior update
\[
p_h(\mathrm{d}y\mid b)=\int h(\mathrm{d}y\mid x)\,b(\mathrm{d}x),
\qquad
\Phi_h(b,y)(\mathrm{d}x)
=\frac{h(\mathrm{d}y\mid x)\,b(\mathrm{d}x)}{p_h(\mathrm{d}y\mid b)}.
\]
The update $\Phi_{h'}(b,y)$ must be defined for $p_h$-a.e.\ $y$, i.e.\
$p_h(\cdot\mid b)$ must be absolutely continuous with respect to
$p_{h'}(\cdot\mid b)$ (and conversely for a symmetric comparison).  The Gaussian
kernel families used in the paper satisfy this.  If a kernel is singular (for
example the HALT atom), first decompose on the common support and treat the
atomic part separately, or use the extended Le Cam deficiency; do not silently
drop the a.e.\ qualifier.  Let $\TV$ and $W_1$ denote total variation and the
$1$-Wasserstein distance, let $|\rho|\le R_{\max}$, and let $B$ bound the value
function used in the one-step decomposition.

\begin{aIIPropMain}
(three-term decomposition: decomposition, not a uniform bound.)
\begin{enumerate}
\item[(a)] \emph{No single uniform bound.}  There is no universal function of a
single $W_1$ distance, and no universal function of a single Le Cam
deficiency, that bounds
$|V_h^\pi-V_{h'}^\pi|$ uniformly over fixed policies $\pi$.  Proposition~A.6
gives two experiments with zero symmetric Le Cam deficiency and a fixed policy
whose value gap is positive; Proposition~A.7 gives two observation laws with
$W_1=q\epsilon$ and a fixed threshold policy whose value gap is $qR_{\max}$, so
no finite Lipschitz constant in $W_1$ exists.
\item[(b)] \emph{Conditional positive decomposition.}  Under the regularity
conditions (R1)--(R3) below, the one-step value gap splits into an estimation
term and an information term (Lemma~A.2).  A fixed observable policy
satisfies the mixed bound \eqref{eq:a2-3.4.2}.  The
first term is controlled by a belief $W_1$ distance; the second by a predictive
distribution TV distance.  The two terms are not interchangeable.
\item[(c)] \emph{Boundary diagnostic.}  For a belief-threshold policy
$\pi(b)$, the policy boundary/action-flip term is a necessary additional term;
for the paper's observable rebalancing policies $\pi(y)$
the TV information term already averages the action flips, so the boundary term
is a diagnostic refinement, not a condition for the bound.  The state-level TV
bound of Proposition~A.1 is a worst-case sensitivity; in the freeze regime it
can be vacuous.
\end{enumerate}
\end{aIIPropMain}

\begin{proof}
Part (a) is proved by Proposition~A.6 and Proposition~A.7.  Part (b) is
Lemma~A.2 and Corollary~A.3.  Part (c) is Proposition~A.8 together with the
policy-class discussion after Corollary~A.3.
\end{proof}

\paragraph{State-level TV sensitivity.}
The following state-level TV sensitivity is a worst-case bound, not the main
guarantee, and it can be empty.

\begin{aIIPropTV}
Assume A1 and that two POMDPs share $(X,A,T,\rho,\nu,\pi)$ and differ only in
the observation kernel, $h$ versus $h'$.  Let
\[
\kappa_{\mathrm{obs}}
:=\sup_{r,x}\TV\bigl(h_r(\cdot\mid x),h'_r(\cdot\mid x)\bigr),
\]
where $x$ ranges over the full latent state.  Then, for finite horizon $H$,
\[
|V_h^\pi-V_{h'}^\pi|
\le R_{\max}(H+1)(H+2)\,\kappa_{\mathrm{obs}}
\le 2R_{\max}(H+1)^2\,\kappa_{\mathrm{obs}},
\]
and for $\gamma<1$,
\[
|V_h^\pi-V_{h'}^\pi|
\le \frac{2R_{\max}\,\kappa_{\mathrm{obs}}}{(1-\gamma)^2}.
\]
\end{aIIPropTV}

\begin{proof}
Couple the two processes as follows.  Use the same initial state
$X_0\sim\nu$ and the same transition random numbers, so that, as long as the
actions agree, the states agree.  At step $t$, if the observation histories up
to $t-1$ have not diverged, use a maximal coupling so that
\[
\Pr(P_t\ne P'_t\mid X_t,\text{ no divergence})
=\TV\bigl(h(\cdot\mid X_t),h'(\cdot\mid X_t)\bigr)
\le\kappa_{\mathrm{obs}}.
\]
If the histories have already diverged, let the processes evolve independently
(or use any coupling).  The policy $\pi$ is the same.  Before divergence the
actions, states, and rewards agree; after an observation differs, the histories
diverge and the reward difference is at most $2R_{\max}$.

Let $D_t=\mathbf{1}\{\text{the history through }t\text{ has diverged}\}$;
$D_t$ covers the $t+1$ observations $P_0,\dots,P_t$.  By the union bound,
\[
\Pr(D_t)\le 1-(1-\kappa_{\mathrm{obs}})^{t+1}
\le (t+1)\kappa_{\mathrm{obs}}.
\]
Therefore
\[
|\mathbb E[R_t]-\mathbb E[R'_t]|
\le \mathbb E\bigl[|R_t-R'_t|\,D_t\bigr]
\le 2R_{\max}\Pr(D_t)
\le 2R_{\max}(t+1)\kappa_{\mathrm{obs}}.
\]
Summing over $t=0,\dots,H$ gives
\[
|V_h^\pi-V_{h'}^\pi|
\le 2R_{\max}\kappa_{\mathrm{obs}}\sum_{t=0}^{H}(t+1)
=R_{\max}(H+1)(H+2)\kappa_{\mathrm{obs}}.
\]
For the discounted case, use
$\sum_{t\ge0}\gamma^t(t+1)=1/(1-\gamma)^2$.
\end{proof}

\paragraph{Caveats on the state-level bound.}
\begin{enumerate}
\item \emph{State-level TV is required.}  The coupling compares
$h(\cdot\mid X_t)$ with $h'(\cdot\mid X_t)$ at the same latent state.  A
belief-mixed quantity such as $\sup_b\TV(h(\cdot\mid b),h'(\cdot\mid b))$ does
not control $\TV(h(\cdot\mid b_t),h'(\cdot\mid b'_t))$ when the beliefs
$b_t,b'_t$ differ; reconnecting it needs an additional Lipschitz condition of
$h$ in $b$.  The belief $W_1$ route is a different object and must not be mixed
with state-level TV.
\item \emph{Constants and filter stability.}  The union bound does not use
filter forgetting or mixing.  Under a belief-mixing assumption (exponential
decay of the initial-condition effect), the divergence probability can be
summed to a constant and the discounted bound improves to
$O(R_{\max}\kappa_{\mathrm{obs}}/(1-\gamma))$; without it,
$1/(1-\gamma)^2$ is the natural constant of this argument.  A skeleton bound
$R_{\max}\kappa_{\mathrm{obs}}/(1-\gamma)$ omits this factor and must not be
quoted without the mixing hypothesis.
\item \emph{The bound can be empty.}  In the synthetic environment the OPEN
kernel concentrates on $S_t$ and the WEEKEND kernel on $S_{\mathrm{last\,open}}$;
when $|S_t-S_{\mathrm{last\,open}}|$ is much larger than the kernel width,
$\TV\approx1$ and the bound is $O(R_{\max}H^2)$, with no numerical force.  This
is not conservatism but vacuity.  Proposition~A.1 is therefore a worst-case
sensitivity, not a main coverage guarantee.
\item \emph{The finite-horizon constant is large.}  At $H=336$,
$2R_{\max}(H+1)^2\approx2.27\times10^5R_{\max}$, and the sharper
$R_{\max}(H+1)(H+2)\approx1.14\times10^5R_{\max}$, both far above the trivial
bound $2R_{\max}(H+1)\approx674R_{\max}$.  The bound is informative only when
$\kappa_{\mathrm{obs}}\ll1$; report $\kappa_{\mathrm{obs}}$ together with the
trivial bound.
\item \emph{The role of A7.}  The coupling requires the same $T$ and no
action effect on the $S$ transition.  With market impact, $T$ under $\pi$ differs
from the logging transition and $V_h^\pi$ itself is not identified; the premise
of Proposition~A.1 fails.  The synthetic environment satisfies A7; real data
must be checked.
\end{enumerate}

\paragraph{One-step decomposition.}
The following lemma is the algebraic core of three-term decomposition.

\begin{aIILemOneStep}
Let $V$ be bounded and measurable, $|V|\le B$.  Then
\[
\begin{aligned}
&\mathbb E_{y\sim p_h(\cdot\mid b)}V(\Phi_h(b,y))
-\mathbb E_{y\sim p_{h'}(\cdot\mid b)}V(\Phi_{h'}(b,y)) \\
&\quad=
\underbrace{\mathbb E_{y\sim p_h(\cdot\mid b)}
\bigl[V(\Phi_h(b,y))-V(\Phi_{h'}(b,y))\bigr]}_{\text{estimation term}} \\
&\qquad+
\underbrace{\mathbb E_{y\sim p_h(\cdot\mid b)}V(\Phi_{h'}(b,y))
-\mathbb E_{y\sim p_{h'}(\cdot\mid b)}V(\Phi_{h'}(b,y))}_{\text{information term}}.
\end{aligned}
\]
If $V$ is $L_V$-Lipschitz in $W_1$, the estimation term is at most
\[
L_V\,\mathbb E_{y\sim p_h(\cdot\mid b)}
W_1\bigl(\Phi_h(b,y),\Phi_{h'}(b,y)\bigr),
\]
and the information term is at most
\[
2B\,\TV\bigl(p_h(\cdot\mid b),p_{h'}(\cdot\mid b)\bigr).
\]
\end{aIILemOneStep}

\begin{proof}
Add and subtract
\[
\mathbb E_{y\sim p_h(\cdot\mid b)}V(\Phi_{h'}(b,y)).
\]
For the estimation term, use the Kantorovich--Rubinstein duality: the difference of integrals of an
$L_V$-Lipschitz function against two probability measures is at most $L_V$ times
their $W_1$ distance.  For the information term, two probability measures give
expectations of the same bounded function $V(\Phi_{h'}(b,\cdot))$ that differ by
at most $2B$ times their TV distance.
\end{proof}

\paragraph{Mixed bound for observable policies.}
Let $T_a$ be the transition kernel under action $a$, and let $bT_a$ be its
pushforward.  For a fixed policy $\pi(a\mid H^o)$ that depends only on the
current observation $y$ (the paper's observable rebalancing rules; if the policy
depends on a longer history, enlarge $b$ to include the policy memory), define
\[
(T_hV)(b)
=\mathbb E_{y\sim p_h(\cdot\mid b)}
\bigl[\rho(b,\pi(y))+\gamma V(\Phi_h(b,y)T_{\pi(y)})\bigr],
\qquad
\rho(b,a)=\mathbb E[\rho(X,a,X')\mid b,a],
\]
where $\rho(b,a)$ is the one-step expected reward given the belief and action
(integrate the transition kernel when the reward depends on $X'$), and
$\Phi_h(b,y)T_{\pi(y)}$ is the \emph{predictive belief}: the current posterior
pushed forward through the transition, not the current posterior itself.

\begin{aIICorMixed}
Assume (R1)--(R3) below.  Then, for $\gamma<1$,
\begin{equation}\label{eq:a2-3.4.1}
\|V_h^\pi-V_{h'}^\pi\|_\infty
\le\frac{1}{1-\gamma}
\|T_hV_{h'}^\pi-T_{h'}V_{h'}^\pi\|_\infty.
\aIIfixedtag{3.4.1}
\end{equation}
If $V_{h'}^\pi$ is $L_V$-Lipschitz in $W_1$ and $|V_{h'}^\pi|\le B$, and $T_a$
is uniformly $L_T$-Lipschitz in $W_1$ for all actions $a$, then
\begin{equation}\label{eq:a2-3.4.2}
\begin{aligned}
\|V_h^\pi-V_{h'}^\pi\|_\infty
\le\frac{1}{1-\gamma}\Bigl[
&\gamma L_VL_T\sup_b\mathbb E_{y\sim p_h(\cdot\mid b)}
W_1\bigl(\Phi_h(b,y),\Phi_{h'}(b,y)\bigr) \\
&+2(R_{\max}+\gamma B)\sup_b
\TV\bigl(p_h(\cdot\mid b),p_{h'}(\cdot\mid b)\bigr)\Bigr].
\end{aligned}
\aIIfixedtag{3.4.2}
\end{equation}
For finite horizon, replace $1/(1-\gamma)$ by the horizon sum (order $H+1$ when
$\gamma=1$).
\end{aIICorMixed}

\begin{proof}
Write
$V_h-V_{h'}=(T_hV_h-T_hV_{h'})+(T_hV_{h'}-T_{h'}V_{h'})$.  Taking the sup norm
and using the $\gamma$-contraction of $T_h$ gives \eqref{eq:a2-3.4.1}.  For the
second term, fix $b$ and write
$\mathbb E_{p_h}[g_h(y)]-\mathbb E_{p_{h'}}[g_{h'}(y)]$ with
\[
g_h(y)=\rho(b,\pi(y))+\gamma V(\Phi_h(b,y)T_{\pi(y)}).
\]
Add and subtract $\mathbb E_{p_h}[g_{h'}(y)]$.  The first segment is
$\gamma\mathbb E_{p_h}[V(\Phi_hT)-V(\Phi_{h'}T)]$, controlled by
$L_VL_T$ and $W_1$.  The second segment is the difference of the expectations of
the same bounded function $g_{h'}$ under two probability measures, with
$|g_{h'}|\le R_{\max}+\gamma B$, hence controlled by
$2(R_{\max}+\gamma B)\TV$.  Taking the supremum over $b$ gives
\eqref{eq:a2-3.4.2}.
\end{proof}

\paragraph{Regularity conditions (R1)--(R3).}
\begin{enumerate}
\item[(R1)] $V_{h'}^\pi$ is $L_V$-Lipschitz in $W_1$ and $|V_{h'}^\pi|\le B$.
\item[(R2)] $T_a$ is uniformly $L_T$-Lipschitz in $W_1$ over actions $a$.
\item[(R3)] $\gamma L_\Phi L_T<1$, where $L_\Phi$ is a uniform $W_1$-Lipschitz
constant of $b\mapsto\Phi_h(b,y)$ in $b$.  A sufficient set of conditions is:
$\rho(b,a)$ is uniformly $L_\rho$-Lipschitz in $b$; $b\mapsto p_h(\cdot\mid b)$
is $L_p$-Lipschitz in TV; $b\mapsto\Phi_h(b,y)$ is uniformly $L_\Phi$-Lipschitz
in $W_1$; and $T_a$ is uniformly $L_T$-Lipschitz in $W_1$.  Then the Bellman
operator maps an $L$-Lipschitz function to an $L'$-Lipschitz function with
\[
L'\le L_\rho+\gamma L L_\Phi L_T+2(R_{\max}+\gamma B)L_p.
\]
If $\gamma L_\Phi L_T<1$, a sufficiently large $L$ satisfies $L'\le L$ and the
fixed point has finite $L_V$.  If $\gamma L_\Phi L_T\ge1$, this recursion does
not guarantee a finite $L_V$; this is exactly the mixing/stability condition
behind the improved discounted constant in the state-level bound, and it cannot
be assumed for free.
\end{enumerate}

\paragraph{Policy classes.}
The bound \eqref{eq:a2-3.4.2} applies to an \emph{observable} policy
$\pi(y)$, even when $\pi$ is a threshold: the threshold discontinuity is in
$y$, and $y$ is integrated out by $p_h(\cdot\mid b)$, so $b\mapsto V_h^\pi(b)$
can still have a finite $L_V$ under kernel regularity.  If the policy is a
\emph{belief} policy $\pi(b)$ (a threshold on the posterior mean or
probability), then $b\mapsto\pi(b)$ is discontinuous, $V_h^\pi(b)$ jumps at the
decision boundary, and $L_V=\infty$; \eqref{eq:a2-3.4.2} does not apply and the
boundary term of Proposition~A.8 must be reported separately.  The two policy
classes must not be conflated.

\paragraph{Reading the two terms.}
The $W_1$ estimation term measures how much the two kernels change the
posterior under the same observation: an estimation effect.  The TV information
term measures how much the two kernels change the predictive law of the next
observation under the same belief: an information effect.  For an observable
policy, the information term also averages out action flips, because the same
$y$ gives the same action and the flip comes only from the change in the law of
$y$.  The two objects are distinct: $W_1(b,b')=0$ can coexist with an
information term of size $2(R_{\max}+\gamma B)$ (same belief, different future
information), so a $W_1$-only bound must miss the information effect; conversely
$\TV(h,h')\approx1$ can coexist with a small $W_1$ estimation term (the prior
width in a freeze period is finite), so a TV-only bound can be vacuous.  These
cases show why the estimation and information terms must be reported separately.

\paragraph{Le Cam: optimal values, not fixed policies.}
\begin{aIIDefLeCam}
For two experiments $E=(P_\theta)$ and $F=(Q_\theta)$ with the same parameter
space $\Theta$, define
\[
\delta(E,F):=\inf_{K}\sup_\theta\TV(KP_\theta,Q_\theta),
\qquad
\delta_{\mathrm{sym}}(E,F):=\max\bigl(\delta(E,F),\delta(F,E)\bigr),
\]
where $K$ ranges over Markov kernels $Y\to Y'$ (randomization is allowed; with
deterministic maps the ``simulate by adding noise'' example fails).
\end{aIIDefLeCam}

\begin{aIILemLeCam}[Bayes optimal values]
For the same prior and a loss $\ell(\theta,a)\in[0,1]$,
\[
V_E^*-V_F^*\le\delta(E,F),
\qquad
V_F^*-V_E^*\le\delta(F,E),
\]
hence $|V_E^*-V_F^*|\le\delta_{\mathrm{sym}}(E,F)$, where $V_E^*$ is the Bayes
optimal value (the sign is irrelevant for the absolute bound).
\end{aIILemLeCam}

\begin{proof}
For any decision rule $\varphi$ under $F$, $\varphi\circ K$ is a rule under
$E$, and the pointwise risk difference is at most $\TV(KP_\theta,Q_\theta)$.
Taking $\inf_\varphi$ and then $\inf_K$ gives the first inequality; exchanging
$E$ and $F$ gives the second.
\end{proof}

For a finite-horizon POMDP, treat the whole observation history as the
experiment output and the policy as the decision rule.  A1 gives
$|R_t|\le R_{\max}$, so the total reward lies in an interval of length
$2(H+1)R_{\max}$; normalize it linearly to $[0,1]$, apply Lemma~A.5, and scale
back:
\begin{equation}\label{eq:a2-3.4.3}
|V_{E_h}^*-V_{E_{h'}}^*|
\le 2(H+1)R_{\max}\,\delta_{\mathrm{sym}}(E_h,E_{h'}).
\aIIfixedtag{3.4.3}
\end{equation}
This requires A7: otherwise actions change future observations, the experiment
is not exogenous, and the Le Cam deficiency cannot be applied to the whole
trajectory.  The position and cash are determined by the action history but can
be reconstructed deterministically from the observation history; they are part
of the decision rule, not additional experiment noise.

The target $V^\pi$ is a fixed-policy value, not an optimal value.  A Le Cam-only
route therefore needs a regret term:
\begin{equation}\label{eq:a2-3.4.4}
|V_h^\pi-V_{h'}^\pi|
\le 2(H+1)R_{\max}\delta_{\mathrm{sym}}(E_h,E_{h'})
+\mathrm{Reg}_h(\pi)+\mathrm{Reg}_{h'}(\pi),
\qquad
\mathrm{Reg}_h(\pi):=V_h^*-V_h^\pi\ge0.
\aIIfixedtag{3.4.4}
\end{equation}

\begin{aIIPropLeCam}[Relabeling counterexample]
Take $\Theta=\{0,1\}$ with a uniform prior and the $0$--$1$ loss
$\ell(\theta,a)=\mathbf{1}\{a\ne\theta\}$.  Experiment $E$: $Y=\theta$;
experiment $F$: $Y=1-\theta$.  Then $\delta(E,F)=\delta(F,E)=0$ (the kernel
$K(y)=1-y$ converts one experiment exactly into the other).  Take the fixed
policy $\pi(y)=y$: its risk is $0$ under $E$ and $1$ under $F$.  The value gap
is therefore $1$ (times $R_{\max}$), while $\delta=0$.  Hence no fixed-policy
value-error bound can depend on $\delta$ alone; the regret term cannot be
dropped.
\end{aIIPropLeCam}

\begin{proof}
The kernel $K(y)=1-y$ is a deterministic Markov kernel and gives
$KP_\theta=Q_\theta$ and $KQ_\theta=P_\theta$, so both deficiencies vanish.  The
fixed policy $\pi(y)=y$ has risk $0$ under $E$ and risk $1$ under $F$, so the
value gap is $1$ after scaling by $R_{\max}$.
\end{proof}

\paragraph{Le Cam asymmetry.}
Le Cam deficiency is asymmetric.  For example, let $E$ be noiseless and let
$F=E+\text{independent noise}$ (so $F$ is less informative).  Then
$\delta(E,F)=0$ (one can simulate $F$ from $E$ by adding noise), but
$\delta(F,E)=1$ (one cannot denoise).  The information value lies in the
$\delta(F,E)$ direction, so the symmetric deficiency is $1$ and
\eqref{eq:a2-3.4.4} is vacuous under information loss.  Le Cam gives a nonempty
optimal-value bound only when the two experiments are mutually simulable
(equal information, e.g.\ relabeling or two-way noise); it cannot rescue the
freeze-period ``informative versus uninformative'' gap.

\paragraph{Why observed $W_1$ cannot replace TV.}
\begin{aIIPropWone}[Threshold counterexample]
Take a one-dimensional observation $Y$ and two observation laws $p_h,p_{h'}$:
$p_h$ has an atom of mass $q>0$ at $0$, and $p_{h'}$ moves that atom to
$\epsilon>0$ (or the reverse), leaving the rest unchanged.  Then
$W_1(p_h,p_{h'})=q\epsilon$ and $\TV(p_h,p_{h'})=q$.  The policy
$\pi(y)=\mathbf{1}\{y>0\}$ flips its action with probability at least $q$.
Choose rewards such that action $1$ in state $0$ is worse than action $0$ by
$R_{\max}$.  Then
\[
|V_h^\pi-V_{h'}^\pi|\ge qR_{\max},
\qquad
W_1(p_h,p_{h'})=q\epsilon,
\qquad
\TV(p_h,p_{h'})=q.
\]
For any $L$, fix $q$ and take $\epsilon<R_{\max}/L$; this violates
$|\Delta V|\le L\,W_1$.  Thus observed-law $W_1$ cannot give a finite Lipschitz
bound for threshold policies, while TV charges the flip with probability $q$
and is of the same order as $qR_{\max}$.
\end{aIIPropWone}

\begin{proof}
The atom displacement gives $W_1=q\epsilon$ and $\TV=q$.  The threshold policy
flips on the displaced atom, an event of probability at least $q$; the reward
loss on that event is $R_{\max}$, so the value gap is at least $qR_{\max}$.  For
fixed $q$, let $\epsilon\downarrow0$ with $L$ fixed; then
$L\,W_1=Lq\epsilon<qR_{\max}$ for $\epsilon<R_{\max}/L$, contradicting the
claimed Lipschitz bound.
\end{proof}

\paragraph{Relation to the mixed bound.}
Proposition~A.7 is not a counterexample to \eqref{eq:a2-3.4.2}: the information
term there uses TV, not $W_1$.  It shows two things.  First, if the information
term is mistakenly written with observed $W_1$, threshold policies defeat any
finite constant.  Second, for observable policies the action flips are already
averaged by the TV term, so a boundary/margin term is a sharper diagnostic, not
an independent bound.  Only for a belief-threshold policy $\pi(b)$ does the flip
occur in the posterior, where the TV predictive term no longer contains it; then
the boundary term must be reported separately (Proposition~A.8).

\paragraph{HALT atoms.}
The counterexample lands directly on the paper's thresholds
($|p/\text{anchor}-1|=0.02$ or $0.01$): the HALT frozen price is an atom (the
pool price does not move during a halt), giving a realistic $q>0$.
In belief space, the same phenomenon is $b=\delta_0$ versus
$b_q=(1-q)\delta_0+q\delta_1$ with $W_1=q$; for a threshold policy $\pi(b)$, the
action flips at the boundary and $V^\pi(b)$ is discontinuous, so $L_V=\infty$.

\begin{aIIPropBoundary}[Boundary-margin condition; sketch]
For a belief policy $\pi(b)$ with decision boundary $\partial\pi$, define the
margin $m(b)=\operatorname{dist}(b,\partial\pi)$ in $W_1$.  If $m(b)>\delta$,
then any perturbation of $W_1$ radius at most $\delta$ leaves the action
unchanged.  Hence the action-flip probability is at most
$\Pr(m(b_t)\le\delta)$, and the value difference can be written as
\[
|V_h^\pi-V_{h'}^\pi|
\le 2R_{\max}\sum_t\gamma^t\Pr\bigl(m(b_t)\le\delta\bigr)
+\text{(posterior-difference term outside the boundary)}.
\]
The probability $\Pr(m(b_t)\le\delta)$ can be estimated directly by simulation.
This is a policy-aware, data-dependent sensitivity, not a prior uniform bound.
For the paper's observable rebalancing policies, this term is not
necessary for \eqref{eq:a2-3.4.2}, because the TV term already contains the
action flips; the value-error track should still report the flip frequency and
margin as a diagnostic refinement of the information term.
\end{aIIPropBoundary}

\begin{proofsketch}
If $m(b)>\delta$, every belief within $W_1$ distance $\delta$ stays on the same
side of the decision boundary, so the action is unchanged.  The only possible
flips are therefore on the event $m(b_t)\le\delta$; on that event the reward
difference is at most $2R_{\max}$, and away from it the remaining discrepancy is
the posterior-difference term.  Summing over $t$ with the discount gives the
displayed bound.  The probability is estimable because $b_t$ and the boundary
are computable from the fitted filter and the known policy.
\end{proofsketch}

\paragraph{Implications for the paper.}
\begin{enumerate}
\item Report the three terms separately: a $W_1$ estimation term, a predictive
TV information term, and a policy boundary/action-flip diagnostic.  Do not
compress them into a single $\kappa_{\mathrm{obs}}$.
\item Do not claim a uniform value-error bound for a fixed policy from
TV, $W_1$, or Le Cam alone.  Proposition~A.1 is a state-level TV worst-case
sensitivity; it is not a main guarantee.
\item The main guarantees remain Theorem~1 (identification and plug-in
consistency) and Theorem~2 (coverage).  The value-error track is an estimable
sensitivity analysis.
\item A conditional positive result exists when $L_V<\infty$ (observable policy
plus kernel regularity and filter contraction) and $W_1$ is controlled: then
\eqref{eq:a2-3.4.2} gives a finite bound.  The paper's observable
rebalancing policies can enter this framework, but in the freeze
regime $\TV\approx1$, so the bound has no numerical force.  Belief-threshold
policies do not satisfy $L_V<\infty$ and need Proposition~A.8.
\item In the freeze regime, the informative-versus-uninformative comparison is
the worst case: OPEN versus WEEKEND has state-level $\TV\approx1$ when $S_t$ is
far from the stale anchor; the belief $W_1$ is finite but the information term
still sees $\TV\approx1$; and the Le Cam deficiency from the uninformative
kernel to the informative one is $1$, so \eqref{eq:a2-3.4.4} is vacuous.  The only
informative guarantee in that regime is the coverage track with
$\kappa_{\mathrm{conf}}$, not a value-error bound.
\end{enumerate}

\let\theHequation\aIItheHequation
